% The Collected Mathematical Papers of Arthur Cayley, Vol. XI
% Transcription of PDF pages 51-100 (printed pages 27-76)
% Papers 712-728

\documentclass[12pt,a4paper,oneside]{book}

% ============================================================
% PREAMBLE
% ============================================================
\usepackage[utf8]{inputenc}
\usepackage[T1]{fontenc}
\usepackage{lmodern}
\usepackage[english]{babel}
\usepackage{amsmath,amssymb,amsthm}
\usepackage{mathtools}
\usepackage{mathrsfs}
\usepackage{geometry}
\usepackage{fancyhdr}
\usepackage{titlesec}
\usepackage{setspace}
\usepackage{array}
\usepackage{booktabs}
\usepackage{longtable}
\usepackage{enumitem}
\usepackage{microtype}

% Geometry
\geometry{
  paperwidth=6.5in,
  paperheight=9.5in,
  top=1in,
  bottom=1in,
  inner=1in,
  outer=0.8in,
}

% Headers and footers
\pagestyle{fancy}
\fancyhf{}
\fancyhead[LE,RO]{\thepage}
\renewcommand{\headrulewidth}{0pt}

% Section formatting - paper numbers
\titleformat{\section}
  {\centering\large\scshape}
  {}{0em}{}
\titlespacing*{\section}{0pt}{2em}{1em}

\titleformat{\subsection}
  {\centering\normalsize\scshape}
  {}{0em}{}
\titlespacing*{\subsection}{0pt}{1em}{0.5em}

\newcommand{\paperref}[1]{\par\noindent\hfill[#1]\hfill\null\par}
\newcommand{\source}[1]{\par\noindent\centering[#1]\par\medskip}

\theoremstyle{plain}
\newtheorem{theorem}{Theorem}
\newtheorem{lemma}{Lemma}
\newtheorem{corollary}{Corollary}

\theoremstyle{definition}
\newtheorem{definition}{Definition}

\theoremstyle{remark}
\newtheorem*{remark}{Remark}

\newcommand{\ph}{\varphi}
\newcommand{\dprime}{\prime\prime}
\newcommand{\tprime}{\prime\prime\prime}
\DeclareMathOperator{\const}{const.}

\title{\vspace{-2em}THE COLLECTED MATHEMATICAL PAPERS\\Vol.~XI --- Pages 27--76}
\author{}
\date{}

\usepackage{graphicx}
\emergencystretch=3em
\tolerance=600
\begin{document}
\setcounter{page}{27}

% ============================================================
% PAPER 712: A PARTIAL DIFFERENTIAL EQUATION CONNECTED
%             WITH THE SIMPLEST CASE OF ABEL'S THEOREM
% PDF page 51 (printed 27)
% ============================================================
\section*{712.}
\subsection*{A PARTIAL DIFFERENTIAL EQUATION CONNECTED WITH THE\\ SIMPLEST CASE OF ABEL'S THEOREM.}

\begin{center}
\small [From the \textit{Report of the British Association for the Advancement of Science}, (1881), pp.~534, 535.]
\end{center}

\medskip

\textsc{Consider} a given cubic curve cut by a line in the points $(x,\ p)$, $(x_1,\ p_1)$, $(x_2,\ p_2)$; using the first and second points as elements, these determine uniquely the third point. Analytically, the equation of the curve determines $p$ as a function of $x$, and $p_2$ as a function of $x_1$; writing in the equation
\[
c_2=\lambda c_1+(1-\lambda)c,\quad y_2=\lambda y_1+(1-\lambda)y_0,
\]
we have $c$, by a simple equation, and thence $p_2$; viz.\ $x_2$ is found as a function of $x$, $x_1$, and of the nine constants of the equation. Hence having the derived equations (in regard to $x$, $x_1$) of the first, second, and third orders, we have $(1+2+3+4=)10$ equations from which to eliminate the 6 quantities $x_2$, considered as a function of $x$ and $x_1$, thus involves a partial differential equation of the third order, independent of the particular cubic curve.

To obtain this equation it is only necessary to observe that we have, by Abel's theorem,
\[
\frac{dx}{X}+\frac{dx_1}{X_1}+\frac{dx_2}{X_2}=0,
\]
where $X$, $X_1$, $X_2$ is a given function of $x$, and $p$, that is, of $x$; $X_1$ and $X_2$ are the like functions of $x_1$ and $x_2$ respectively. Hence, considering $x_2$ as a function of $x$ and $x_1$, we have
\[
\frac{dx_2}{dx}=-\frac{X_2}{X},\quad \frac{dx_2}{dx_1}=-\frac{X_2}{X_1},
\]

\paperref{\textit{C. XI.}\hfill 4--2}

\newpage
% ============================================================
% PDF page 52 (printed 28)
% ============================================================
and consequently
\[
\frac{dx_2}{dx_1}\div\frac{dx_2}{dx}=\frac{X}{X_1},
\]
where $X$, $X_1$ are functions of $x_2$, $x_1$ respectively; hence taking the logarithm and differentiating successively with regard to $x_1$ and $x_2$, we have
\[
\frac{d}{dx_2}\frac{d}{dx_1}\log\left(\frac{dx_2}{dx_1}\div\frac{dx_2}{dx}\right)=0,
\]
which is the required partial differential equation of the third order.

This differential equation has a simple geometrical significance. Consider three consecutive positions of the line meeting the cubic curve in the points 1, 2, 3; $1'$, $2'$, $3'$; $1''$, $2''$, $3''$ respectively; \textit{qu\ae{}} equation of the third order, the equation should in effect determine $3''$ by means of the other points. And, in fact, the three positions of the line constitute a cubic curve; the nine points are thus the intersections of two cubic curves, or, say, they are an ``annexed'' of points; any eight of the points thus determine uniquely the ninth point.

\vspace{2em}
\noindent\dotfill

% ============================================================
% PAPER 713: ADDITION TO MR ROWE'S MEMOIR ON ABEL'S THEOREM
% PDF page 53 (printed 29)
% ============================================================
\section*{713.}
\subsection*{ADDITION TO MR ROWE'S MEMOIR ON ABEL'S THEOREM.}

\begin{center}
\small [From the \textit{Philosophical Transactions of the Royal Society of London}, vol.~CLXII., Part~III. (1881), pp.~715---758. Received May 27,---Read June 16, 1880.]
\end{center}

\medskip

\textsc{In} Abel's general theorem $y$ is an irrational function of $x$ determined by an equation $\chi(y)=0$, or say $\chi(x,\ y)=0$; if the order $n$ as regards $y$, and that is also by him that the sum of any number of the integrals considered may be reduced to a sum of $\rho$ integrals; where $\rho$ is a determinate number depending only on the form of the equation $\chi(x,\ y)=0$, and given in his equation (42), [\textit{(Euvres Compl\`etes}, (1881), t.~I., p.~157]; viz.\ if, solving the equation so as to obtain from it developments of $y$ in descending series of powers of $x$, we have:
\begin{align*}
v_\sigma\rho_\nu,\ \text{series each of the form } y&=x^{\sigma'_\nu}+\dots,\\
v_\sigma\rho_2,\ &\qquad y=x^{\sigma'_2}+\dots,\\
\vdots\qquad&\qquad\vdots\\
v_\sigma\rho_n,\ &\qquad y=x^{\sigma'_n}+\dots,
\end{align*}

\footnotetext{The second power of a term coefficient: this being a relation $y=x^p+\dots$, representing the $y$, different values of $y$ obtained by giving to the radical $x^p$ and to its $p$ values, and the corresponding values in the radicals which enter into representing $y$, values of $y$. It is assumed that the equation is such that the meaning is to make each representing $y$, $y_2$. In each of the equations the meaning of $y$ is taken positive in the development of $y$ in descending powers of $x$.}

\newpage
% ============================================================
% PDF page 54 (printed 30)
% ============================================================
(so that $\sigma=v_1\rho_1+v_2\rho_2+\dots+v_n\rho_n$) then $y$ is a determinate function of $x_1$, $x_2$, $x_3$, $x_4$, $x_5$, $x_6$, $\dots$, $x_n$, say $\rho$.

We have so expressed Abel's $\rho$ in the following form, viz.\ assuming
\[
y=\Sigma a_1 \alpha_1 \nu_1+\Sigma a_2\nu_2\dots+\Sigma a_n\nu_n-\frac{1}{2}\Sigma a-\frac{1}{2}n+1,
\]
or, what is the same thing, for $\rho$ writing this $\Sigma a_1$,
\[
y=\Sigma a_1\nu_1+\Sigma a_2\nu_2+\Sigma a_3\nu_3+\dots+\Sigma a_n\nu_n-\frac{1}{2}\Sigma a-\frac{1}{2}n+1,
\]
where in the first sum $\nu_1$ have each of them the values 1, 2, $\dots$, $\rho_1$ subject to the condition $i+\nu_1\sigma'_1$; in each of the other sums $\nu_i$ and $\rho_n$ are interrelated by the suffix, which has the values 1, 2, etc.

It is a leading result in Riemann's theory of the Abelian integrals that $y$ is the deficiency (Geschlecht) of the curve represented by the equation $\chi(x,\ y)=0$; and it must consequently be demonstrable \textit{a posteriori} that the foregoing expression for $y$ is in fact $=$~deficiency of curve $\chi(x,\ y)=0$. I propose to verify this by means of the formulae given in my paper ``On the Higher Singularities of a Plane Curve,'' \textit{Quart.\ Math.\ Journ.}, vol.~VII., (1866), pp.~212---222, [374].

It is necessary to distinguish between the values of $\frac{m}{n}$ which are $>,=,$ and $<1$; and to fix the ideas I assume $l=7$, and
\begin{gather*}
\frac{m_1}{n_1},\ \frac{m_2}{n_2},\ \frac{m_3}{n_3},\ \text{each}>1,\\
\frac{m_4}{n_4}=1;\ \text{say } m_4=n_4=n,\text{ and } n_4=n;\\
\frac{m_5}{n_5},\ \frac{m_6}{n_6},\ \frac{m_7}{n_7},\ \text{each}<1,
\end{gather*}
but it will be easily seen that the reasoning is quite general. And now $\Sigma$ to denote a sum in regard to the first set of suffixes 1, 2, 3, and $\Sigma'$ to denote a sum in regard to the second set of suffixes 5, 6, 7. The foregoing value of $\rho$ is thus
\[
y=\Sigma a\nu+4+\Sigma'a\nu.
\]
Introducing a third coordinate $z$ for homogeneity, the equation $\chi(x,\ y)=0$ of the curve will be
\[
0=\left[y^{p-(n-1)},\ y^{n-2}x^a,\ \dots\right]\left(y,\ z\right)^{n-2}+\left(y^{n-2},\ y^{n-3}x^a,\ \dots\right)\left(y,\ z\right)^{n-1}+\dots+\left(y\right)\left(y,\ z\right)^n,
\]
where it is to be observed that $\chi(x,\ y)=0$ is the product of $\nu_i$, different ovaries; these divide themselves into $n_i$

\newpage
% ============================================================
% PDF page 55 (printed 31)
% ============================================================
groups, each a product of $\rho_i$ orates; and in each such product the $\rho_i$ coefficients $A_i$ are in general the $\rho_i$ values of a function containing a radical $a^{1/\rho_i}$ and are thus different from each other; it is to when follows in effect assumed in each step that all the $a_i$, but that all the $x_i$ coefficients, etc.\ are different from each other.

The remarks apply to the other factors. It applies in particular to the factor $\left(y,\ z\right)^{r-1}$, viz.\ it is assumed that the coefficients $A$ in the radical $z$, $y^{n-2}$, $\dots$ are all of them different from each other. These assumptions as to the leading coefficients really imply Abel's assumption that $\frac{m_1}{n_1}$, $\frac{m_2}{n_2}$ are all of them fractions in their least terms, viz.\ that $n=1$, since
\[
1\textit{ rotate however for convenience the equation will be considered, by giving to the leading coefficient (in $z$, $y$,)}\ a\ \textit{form is a fraction in its least terms, $\frac{m_i}{n_i}$ is taken as $\frac{m_i}{n_i}=\frac{n_i+1}{n_i}$.}
\]

In the product of the several infinite series, the terms containing the lowest powers of $z$, $\frac{1}{n_1}$ is a function in its lowest terms, viz.\ since $n_1=1$.

A curve has singularities (singular points) at infinity, that is, on the line $z=0$; vis.\
\textit{First}, a singularity at $(x=0, x'=0)$, where the tangent is $x'=0$, and which, writing for convenience $x=1$, is denoted by the function
\[
\left(z-\nu^{m/n}x^{1-1/n}\right)^{n-1}=\dots
\]
where observe that the expressed factor indicates $r_n$ branches $\left(z-y^{m/n}x^{1-1/n}\right)^{r_n}$, or say $r_n(m_n-p_n)$ partial branches $z=v^{m/n}\dots$, that is, $r_n(m_n-p_n)$ partial branches $z=v^{m/n}\dots$, will in all $r_n(m_n-p_n)+\dots$, where $r_n(m_n-p_n)$ partial branches at the encompassed factors with the suffixes $1'$ and $4$; that is, we can regard the encompassed factors with the suffixes $1'$ and $4$, where in the equation
\[
z-\nu^{m/n}+\dots=0.
\]

\textit{Secondly}, a singularity at $(z=0, y=0)$, where the tangent is $y=0$, and which, writing for convenience $z=1$, is denoted by the function
\[
\left(y-z^{m/n}x^{1-1/n}\right)^{n-1}=\dots
\]

\footnotetext{* This assumption is virtually made by Abel, (1), [c.~p.~157], in considers the equation $z^n=y^n\sqrt[m_n]{f(x)}$ as denoting integers; but where it appears that the indices $a$, $b$, are to be considered of fixed, and as in the values of degrees with respect to $x$. He says that if these be assumed as integers in the expression of $y$ in descending powers of $x$ (the leading coefficient being $z$ in regard to $y$), the singular factor are in $(x,\ y)$ in regard to $y$ in regard to $x$.}

\newpage
% ============================================================
% PDF page 56 (printed 32)
% ============================================================
where observe that the expressed factor indicates $r_5$ branches $\left(y-z^{m/n}x^{1-1/n}\right)^{r_5}$, or say $r_5(n_5-m_5)$ partial branches $y=z^{m/n}\dots$, that is, $r_5(n_5-m_5)$ partial branches $y=z^{m/n}\dots$, with in all $r_5n_5+\dots$, with $r_5n_5+r_6n_6+r_7n_7$ partial branches $y=z^{m/n}\dots$, with the suffixes $5$, $6$, $7$.

\textit{Thirdly}, suppose at $(x=0, y=0)$, where the tangent is $z=0$, and consequently $r_4n_4=0$, partial branches $z=y^0\dots+$ similar terms with respect to $x$, $y$, $z$.

It is convenient to assume $r_4=2$, so that
\begin{align*}
E&=\tfrac{1}{2}(K-1)(K-2)=z,\\
E'&=\tfrac{1}{2}(K'-1)(K'-2)\\
&=\tfrac{1}{2}\Sigma(n-1).
\end{align*}
hence
\[
E+e'=\tfrac{1}{2}(K-1)+\tfrac{1}{2}\Sigma(n-1).
\]
Assuming that there are no other singularities, the deficiency
\[
{}_\rho=\tfrac{1}{2}(K-1)(K-2)-E-c'
\]
is
\[
=\tfrac{1}{2}(K-1)(K-2)-\tfrac{1}{2}KK+\tfrac{1}{2}\Sigma(n-1).
\]
This should be equal to the before-mentioned value of $\rho$; viz.\ we ought to have
\[
\tfrac{1}{2}(K-1)(K-2)-\tfrac{1}{2}KK+\tfrac{1}{2}\Sigma(n-1)=\Sigma a\nu_1+4+\Sigma'a\nu_2-\Sigma a-\Sigma'a-2,
\]
or, as it will be convenient to write it,
\[
M=K-\tfrac{3}{2}K+\tfrac{1}{2}\Sigma(n-1)-1=\Sigma a\nu_1+\Sigma'a\nu_2-\Sigma'a-2,
\]
which is the equation which ought to be satisfied by the values of $N$ and $\Sigma (n-1)$ calculated, according to the method of my paper, for the foregoing singularities of the curve.

We have as before:
\[
K=\Sigma n+\Sigma'n+\theta n.
\]
The terms $\Sigma n+\Sigma'n+\theta n$, written at length, is
{\small
\begin{align*}
&=n_1(\nu_1\rho_1+\nu_2\rho_2+\nu_3\rho_3)+\dots+n_2(\nu_1\rho_1+\nu_2\rho_2+\nu_3\rho_3)\\
&+\dots+n_3(\nu_1\rho_1+\nu_2\rho_2+\nu_3\rho_3)\\
&+\dots+n_4(\nu_1\rho_1+\nu_2\rho_2+\nu_3\rho_3)\\
&+\dots+n_5(\nu_1\rho_1+\nu_2\rho_2+\nu_3\rho_3)\\
&+\dots+n_6(\nu_1\rho_1+\nu_2\rho_2+\nu_3\rho_3)\\
&+\dots+n_7(\nu_1\rho_1+\nu_2\rho_2+\nu_3\rho_3).
\end{align*}}

\newpage
% ============================================================
% PDF page 57 (printed 33)
% ============================================================
which is
\[
=\Sigma\nu_1n_2+\Sigma'n_2(2\rho_2+2'\rho_1)+\Sigma\Sigma'+\Sigma'\Sigma+\Sigma a_1 a_2 a_3,
\]
We have moreover
\[
\Sigma'\rho_1=\Sigma a_1\rho_1+(\theta+\Sigma' a_1)\rho_1=\Sigma a_1+\Sigma a_2+\Sigma'a_1\rho_1,
\]
each branch $\left(z-y^{m/n}x^{1-1/n}\right)^{n-1}$ gives $a=m_1\rho_1$, and the value of $\Sigma(\kappa-1)$ for this singularity is
\[
n_1(m_1-n_1-1)+r_1n_1(m_1-1)=r_1n_1\rho_1-n_1\rho_1.
\]
which is
\[
=\Sigma a\nu_1-\Sigma a\rho_1.
\]
each branch $\left(y-z^{m/n}x^{1-1/n}\right)^{n-1}$ gives $\sigma=m_n\rho_n$, and the value of $\Sigma(\kappa-1)$ for this singularity is
\[
n_4(m_4-n_4-1)+r_4n_4(m_4-1)=r_4n_4\rho_4-n_4\rho_4,
\]
which is
\[
=\Sigma'\nu_2-\Sigma'\nu_4-\Sigma'\nu_1.
\]
For each of the $\theta$ singularities
\[
\left(y-z^{m/n}x^{1-1/n}\right)^{n-1}.
\]
we have $a=n$, and the value of $\Sigma(\kappa-1)$ is $r=\theta(\kappa-1)$; this is $=0$ for the value $\lambda=1$, which is physically attributed to $\lambda$.

According to the theory explained in my paper above referred to, the several singularities not together equivalent to a certain number $K$ of $\kappa^*$ nodes and cusps; viz.\ we have
\[
K=\Sigma n-\Sigma a+\tfrac{1}{2}\Sigma(\kappa-1).
\]
This should be equal to the before-mentioned value of $\rho$; viz.\ we ought to have
\[
(K-1)(K-2)-2(K+\Sigma n+\Sigma'n+\theta n)=2\Sigma a\nu_1+4+2\Sigma a\nu_2-2\Sigma a-2\Sigma'a-4,
\]
or to it will be convenient to write
\[
M-K-\Sigma K+1+\Sigma(n-1)=\Sigma a\nu_1+\Sigma'a\nu_2-\Sigma a-\Sigma'a-2,
\]
which is the equation which ought to be satisfied by the values of $N$ and $\Sigma(n-1)$ calculated, according to the method of my paper, for the foregoing singularities of the curve.

We have as before:
\[
K=\Sigma n+\Sigma'n+\theta n.
\]
The term $\Sigma n+\Sigma'n+\theta n$, written at length, is
\[
=\Sigma\nu_1n_1+\Sigma'\nu_2n_2+\theta n,
\]
\paperref{C.~XI.\hfill 5}

\newpage
% ============================================================
% PDF page 58 (printed 34)
% ============================================================
or, reducing,
\[
M=(\Sigma m)^2-\Sigma m a-\Sigma m \rho_1-\Sigma m \rho_2+(\Sigma' m)^2-\Sigma' m_2 a-\Sigma' m_2 \rho_2
\]
\[
+(\Sigma' m)^2-\Sigma' m\rho_1-\Sigma m a+2\Sigma' m_2 a_1+2\Sigma'\Sigma a_1 a_2;
\]
and it is to be shown that the two lines of this expression are in fact the values of $M$ belonging to the singularities
\[
\left(z-y^{m/n}x^{1-1/n}\right)^{n-1},\ \dots,\ \text{and}\ \left(y-z^{m/n}x^{1-1/n}\right)^{n-1},\ \dots,
\]
respectively. We assume $\lambda=1$, and there is then no singularity $\left(y-z\right)^{n-1}$.

I recall that, considering the several partial branches which meet in a singular point, $M$ denotes the sum of the number of the intersections of each partial branch by every other partial branch: so that for each pair of partial branches the intersections are to be counted twice. Supposing that the tangent is $x=0$, and that, for any two branches we have $z=Az^x+\dots$ ($x$ where $p_x$, $p_x$ are each up to greater than 1), then if $p_x=p_x$, and $A_x=A_x$, $A_x=A_x$ where $A_x, A_x$ are not the partial branches which have here always could be regards the cases about to be considered), then the number of intersections is taken to be $-$; and if $p_x>p_x$, and the number of intersections (taken to be $\equiv p_x$) viz.\ in the case of as a unequal exponents, it is equal to the smaller exponents.

Consider now the singularity $\left(z-y^{m/n}x^{1-1/n}\right)^{n-1},$ and first the intersections of a partial branch $z=y^{m/n}\dots$ by each of the remaining $n_1(m_1-p_1)-1$ partial branches of the same set: the number of intersections with any one of these is $=\frac{m_1}{n_1}\left[n_1(m_1-p_1)-1\right]$. But we obtain this same number from each of the $n_1(m_1-p_1)$ partial branches, and thus the whole number is
\[
n_1\left(m_1-p_1\right)\frac{m_1}{n_1}\left[n_1(m_1-p_1)-1\right]=n_1m_1\left[n_1(m_1-p_1)-1\right].
\]
Taking account of the other sets, and the whole number of intersections is
\[
=\Sigma n^2 a-\Sigma' n a^2.
\]

\newpage
% ============================================================
% PDF page 59 (printed 35)
% ============================================================
Observe now that $\frac{m_1}{n_1}>\frac{m_2}{n_2}$, that is, $\frac{n_1}{n_2}<\frac{m_1}{m_2}$, and that, these being each $<1$, we have $1-\frac{m_2}{n_2}>1-\frac{m_1}{n_1}$, that is, $\frac{n_2-m_2}{n_2}>\frac{n_1-m_1}{n_1}$, and we thus have
\[
\frac{m_1}{n_1}<\frac{m_2}{n_2}, \text{ and } \left(z-y^{m/n}x^{1-1/n}\right)^{n_1(m_1-p_1)},
\]
Considering now the intersections of partial branches of the two sets
\[
\left(z-y^{m/n}x^{1-1/n}\right)^{n_1(m_1-p_1)}\ \text{and}\ \left(z-y^{m/n}x^{1-1/n}\right)^{n_2(m_2-p_2)},
\]
respectively, a partial branch $z-z^{m/n}\dots$ gives with each partial branch of the other set a number $=\frac{m_2}{n_2}$; and in this way taking each partial branch of each set, the number is
\[
n_1(m_1-p_1)\cdot n_2(m_2-p_2)\cdot \frac{m_2}{n_2}=n_2 m_2 n_1(m_1-p_1),
\]
and thus for all the sets the number is
\[
=2\Sigma n_2 m_2 \nu_2 m_2,
\]
where in the first sum the $\Sigma'$ refers to each pair of values of the suffixes. But the intersections are to be taken twice; the number thus is
\[
=2\Sigma m_1 m_2 \nu_1 \nu_2,
\]
Adding the foregoing number
\[
=\Sigma' n^2 a-\Sigma' n a^2;
\]
the whole number for the singularity in question is
\[
=(\Sigma m_1)^2-\Sigma' m_2 \rho_2-\Sigma' m_2 \rho_2 \cdot
\]
Similarly for the singularity $\left(y-z^{m/n}x^{1-1/n}\right)^{n-1}\dots$, taking each set with itself, the number of intersections is
\[
n_5 m_5 \left[n_5(m_5-p_5)-1\right] + n_6 m_6 \left[n_6(m_6-p_6)-1\right]+ n_7 m_7\left[n_7(m_7-p_7)-1\right],
\]
which is
\[
=\Sigma' m_2^2 a-\Sigma' m_2 \rho_2-\Sigma m_1 a.
\]

\paperref{\hfill $5-2$}

\newpage
% ============================================================
% PDF page 60 (printed 36)
% ============================================================
We have here $\frac{m_5}{n_5}>\frac{m_6}{n_6}$, each of these being $<1$, we have $1-\frac{m_5}{n_5}<1-\frac{m_6}{n_6}$, that is, $\frac{n_5-m_5}{n_5}<\frac{n_6-m_6}{n_6}$, and so
\[
\frac{m_5}{n_5}>\frac{m_6}{n_6}, \frac{n_5-m_5}{n_5}<\frac{n_6-m_6}{n_6};
\]

Hence considering the two sets
\[
\left(y-z^{m/n}x^{1-1/n}\right)^{n_5(m_5-p_5)}\ \text{and}\ \left(y-z^{m/n}x^{1-1/n}\right)^{n_6(m_6-p_6)};
\]
a partial branch of the first set given with a partial branch of the second set a number of intersections; and the number then obtained is
\[
n_5(m_5-p_5)\cdot n_6(m_6-p_6)\cdot \frac{m_6}{n_6}=n_6 m_6 n_5(m_5-p_5),
\]
For all the sets the number is
\[
n_5 m_5 \left[n_5 (m_5-p_5)+ n_6(m_6-p_6) \right] - n_5 n_6 (m_6 - m_5)
\]
or taking these from, the number is
\[
2\Sigma'm_2 m_2 \nu_2 \nu_2.
\]
where in the first sum the $\Sigma'$ refers to each pair of suffixes. Adding the foregoing value
\[
\Sigma' m^2 a-\Sigma' m_2 \rho_2-\Sigma m_1 a,
\]
the whole number for the singularity in question is
\[
=(\Sigma' m_2)^2-\Sigma m_1 a-\Sigma' m_2 \rho_2 \cdot
\]
and the proof is thus completed.

Referring to the first-note (ante, p.~31), I remark that the theorem ``$\rho$--deficiency, is absolute, and applies to a curve with any singularities whatever: in a curve which has singularities not taken account of in Abel's theory, the ``perhaps see particulars'' which we have a give as he expresses it, taken account of, viz.\ for that $\Sigma a$ is, in the value of $\rho$ as denoted by Abel's formula; and the actual deficiency will be $=$ Abel's $\rho$, $=$ such diminution, that is, it will be $=$ true values of $\rho$.

\vspace{2em}
\noindent\dotfill

\newpage
% ============================================================
% PAPER 714: VARIOUS NOTES (PDF p61, printed 37)
% ============================================================
\section*{714.}
\subsection*{VARIOUS NOTES.}

\begin{center}
\small [From the \textit{Messenger of Mathematics}, vol.~VII. (1878), pp.~69, 115, 124, 125.]
\end{center}

\medskip

\subsection*{\textit{An Identity.}}

\textsc{The} following remarkable identity is given under a slightly different form by Gauss, \textit{Werke}, t.~III., p.~419.
\begin{align*}
1&=\left(\frac{1}{x}\right)^2x+\left(\frac{1}{1+x}\right)^2x+\left(\frac{1+x}{1+x+x^2}\right)^2x+\dots\\
&=\left[1+\left(\frac{1}{x}\right)^2x+\left(\frac{1}{1+x}\right)^2x+\left(\frac{1+x}{1+x+x^2}\right)^2x+\dots\right]^2.
\end{align*}

\subsection*{\textit{On two related quadric functions.}}

Assume
\begin{align*}
\ph x&=x^2(1-x)+x(1-x^2-\alpha x),\\
\ph y&=y(1-x)+y(x^2-1-\alpha x).
\end{align*}
then
\[
\ph\left(\frac{x^2}{x-a}\right)=\frac{x^2}{(x-a)^2}\ph x,
\]
\[
\ph\left(\frac{x^2}{x-a}\right)=\frac{x^2}{(x-a)^2}\ph y.
\]

In the like of those for a ratio $\frac{p}{1-p}$, then
\[
\ph\left[\frac{ax^2(x-y)}{x^2-y^2(x-a)}\right]=\frac{ax^2(x-y)^2}{(x-a)^2-y^2(x-a)^2}\,\ph(x).
\]

\newpage
% ============================================================
% PDF page 62 (printed 38)
% ============================================================
\subsection*{\textit{A Trigonometrical Identity.}}

\begin{align*}
\cos(b-c)\sin(b+c-a)+\cos a\sin(a+b)\\
+\cos(c-a)\sin(c+a+b)+\cos b\sin(b+c)\\
+\cos(a-b)\sin(a+b+c)+\cos c\sin(c+a)\\
=\cos a\cos b\cos(b+c)+\cos b\cos c\cos(c+a)+\cos c\cos a\cos(a+b)+\cos a\cos b\cos c.
\end{align*}

\subsection*{\textit{Extract from a Letter.}}

``I wish to construct a correspondence such as
\[
(x+iy)^2(x+iy)=X+iY,
\]
or, say, for greater convenience
\[
4(x+iy)^2(x+iy)=X+iY,
\]
viz.\ if
\[
x+iy=\cos\alpha,
\]
then
\[
X+iY=\cos\alpha\beta.
\]

Suppose $\alpha$, in a value of $\alpha$ corresponding to a given value of $X+iY$, then the three values of $x+iy$ are of course $\cos\alpha$, $\cos\left(\alpha+\frac{2\pi}{3}\right)$, but I am afraid that the calculation of $\alpha$, even with red and card tables, would be very laborious. Writing
\[
X+iY=R(\cos\Theta+i\sin\Theta),
\]
the intervals for $\Theta$ might be $3^\circ$, $30^\circ$ or even $15^\circ$; those of $R$, say $0.1$ from $0$ to $2$, and then $0.2$ up to $4$; and 2 places of decimals would be quite sufficient; but even this would probably involve a great mass of calculation.

It has occurred to me that perhaps a geometrical solution might be found for the equation $X+iY=\cos\beta$.''

\hfill\textit{October 21, 1877.}

\vspace{2em}
\noindent\dotfill

\newpage
% ============================================================
% PAPER 715: NOTE ON A SYSTEM OF ALGEBRAICAL EQUATIONS (PDF p63, printed 39)
% ============================================================
\section*{715.}
\subsection*{NOTE ON A SYSTEM OF ALGEBRAICAL EQUATIONS.}

\begin{center}
\small [From the \textit{Messenger of Mathematics}, vol.~VII. (1878), pp.~17, 18.]
\end{center}

\medskip

\textsc{Assume}
\begin{align*}
x+y+z&=P,\\
yz+zx+xy&=Q,\\
xyz&=R.
\end{align*}
\begin{align*}
A&=(nx+y-z)-a^n(nx+P),\\
B&=y(ax+Q)-b^n(my+Q),\\
C&=z(ax+Q)-c^n(nz+P),\\
\Theta&=-mnQ+nP.
\end{align*}
Then
\begin{align*}
(ax+P)B-(my+P)A &=(ax+P)y(ax+Q)-a^n(my+Q)\\
&=(ax+P)(ax+Q)-(my+Q)(my+P)\\
&=ax(my+Q-ay)-yz(my+Q-P)+yz(ax+P)-y(my+Q-Q)\\
&=mny(ax-y)-P(ax-y)\\
&=(ax-y)(mny-P)=(ax-y)\Theta,
\end{align*}
whence, identically,
\begin{align*}
(ax+P)B-(my+P)A&=(ay-y)\Theta,\\
(my+P)C-(nx+P)B&=(by-z)\Theta,\\
(ax+P)A-(nx+P)C&=(cz-x)\Theta.
\end{align*}
Hence any two of the equations $A=0$, $B=0$, $C=0$ imply the third equation.

\newpage
% ============================================================
% PDF page 64 (printed 40)
% ============================================================
We have
\begin{align*}
A&=x\left[(n+1)x+y+z\right]-a^n\left[(n+1)x+y+z\right]\\
&=(x-a^n)\left[(n+1)x+y+z\right];
\end{align*}
and similarly for $B$ and $C$. The three equations therefore are
\begin{align*}
\frac{x}{x^n-a^n}&=\frac{y+z}{(n+1)x^n+(n-1)a^n},\\
\frac{y}{y^n-b^n}&=\frac{x+z}{(n+1)y^n+(n-1)b^n},\\
\frac{z}{z^n-c^n}&=\frac{x+y}{(n+1)z^n+(n-1)c^n},
\end{align*}
and any two of these equations imply the third equation.

\vspace{2em}
\noindent\dotfill

\newpage
% ============================================================
% PAPER 716: AN ILLUSTRATION OF THE THEORY OF \vartheta-FUNCTIONS
% PDF p65, printed 41
% ============================================================
\section*{716.}
\subsection*{AN ILLUSTRATION OF THE THEORY OF THE $\vartheta$-FUNCTIONS.}

\begin{center}
\small [From the \textit{Messenger of Mathematics}, vol.~VII. (1878), pp.~27---32.]
\end{center}

\medskip

\textsc{If} $X$ be a given quartic function of $x$, and if $a$, as for convenience a constant multiple as, be the value of the integral $\int\frac{dx}{\sqrt{X}}$, taken from a given inferior limit to the superior limit $x$; then, conversely, $x$ is expressible as a function of $a$, viz.\ it is expressible in terms of $\vartheta$-functions of $a$, where $\vartheta a$, or say $\vartheta(a,\ \mathfrak{Q})$, is the definition of which, depending, depend in general by definition on the sum of $a$ series of exponentials of $a$, and it is possible from the assumed equation $a=\int\frac{dx}{\sqrt{X}}$, and the definition of $\vartheta a$, to attain by general theory the actual formulae for the determination of $a$ as such a function of $a$.

I propose here to obtain these formula, in the case where $X$ is a product of real factors, in a less suitable manner, viz.\ by reducing the integral $\int\frac{dx}{\sqrt{X}}$ by a linear substitution in the form of an elliptic integral, the object being to obtain for the same in question the actual formulae as a function of $\vartheta$-functions of $u$.

The definition of $\vartheta u$, when the parameter is expressed, $\vartheta(u,\ \mathfrak{Q})$ is
\[
\vartheta u=2\Sigma(-)^p q^{p+\frac{1}{2}(p+1)^2}\sin(p+\tfrac{1}{2})u,
\]
where $u$ and $p$ are negative integer values, zero included, from $-\infty$ to $+\infty$ (that is, from $-2$ to $2$, $0$ to $2$, etc.); the parameter $\mathfrak{Q}$, or $\Omega'$ (of imaginary) no real part, must be positive.

\paperref{C.~XI.\hfill 6}

\newpage
% ============================================================
% PDF page 66 (printed 42)
% ============================================================
Evidently $\vartheta u$ is an even function: $\vartheta(-u)=\vartheta u$. Moreover, it is at once seen that we have
\[
\vartheta(u+\pi)=\vartheta u,\quad \vartheta(u+i\Pi)=-q^{-1}e^{-iu}\vartheta u,
\]
where $q$ is
\[
\vartheta(u+\pi+i\Pi).
\]
where as $u$ and $v$ are positive or negative integers, the product of $\vartheta u$ by an exponential factor, or say simply that it is a multiple of $\vartheta u$.

Writing $u=\frac{\pi}{2}\xi$, we have $\vartheta v=\frac{1}{2}\sqrt{2K}\vartheta\xi$, that is,
\[
\vartheta\left(\tfrac{\pi}{2}\xi\right)=\xi^{nK},
\]
and therefore also
\[
\vartheta(\sin\xi+v)=\xi^{nK}\vartheta v.
\]
The above properties are general, but if $\theta$ be real, then $K, K', q$ being as in Jacobi (consequently $q$ being real, positive, and less than 1, and $K$ and $K'$ real and positive), and assuming $\mathfrak{Q}=\frac{K'}{K},\pi$, we have in the same thing,
\[
q(=e^{-\pi\Omega/\mathfrak{Q}})=e^{-\pi K'/K},
\]
the function $\vartheta v$ is given (\textit{ut} ante) of Jacobi's $\vartheta v$ by the equation $\vartheta u=q^{\frac{1}{4}}\vartheta x$.

That is, if $\Omega$ is the value of the case $u=v$.

We have, as before, the system of three equations
\[
\operatorname{sn} u=\frac{1}{\sqrt{k}}\frac{\vartheta(u,\ k)}{\vartheta_0(u,\ k)},\quad \operatorname{cn} u=\sqrt{\frac{k'}{k}}\,\frac{\vartheta_0(u,\ k)}{\vartheta_0(u,\ k)},\quad \operatorname{dn} u=\sqrt{k'}\,\frac{\vartheta_0(u,\ k)}{\vartheta_0(u,\ k)}.
\]

Consider now the integral
\[
\int_a^b\frac{dx}{\sqrt{(x-\alpha,\ x-\beta,\ x-\gamma,\ x-\delta)}}=\int_a^b\frac{dx}{\sqrt{X}}\text{ suppose,}
\]
where $\alpha$, $\beta$, $\gamma$, $\delta$ are taken to be real, and the four factors regarded as fixed magnitudes, viz.\ it is assumed that $b-a$ is the value of the radical at the inferior limit is positive, and as $b$ is between $\beta$ and $\delta$, $X$ is positive, and \textit{a} fortiori then $\sqrt{X}$ denotes the positive value of the radical; but if $\alpha$ is between $\beta$ and $\delta$, then $X$ is negative, and we must

\newpage
% ============================================================
% PDF page 67 (printed 43)
% ============================================================
that the sign of $\sqrt{X}$ is taken so that $\sqrt{X}$ is equal to a positive multiple of $i$, and this being so the integral is taken from the inferior limit to the superior limit $x$, which we may.

Take a linear function of $y$, such that for
\[
x=a,\ b,\ c,\ d,
\]
\[
y=0,\ 1,\ \tfrac{1}{k^2},\ \infty,\ \text{respectively},
\]
so that, $x$ increasing continuously from $a$ to $d$, $y$ will increase continuously from $0$ to $\infty$. We have
\[
y=\frac{c-a}{c-b}\cdot\frac{x-b}{x-a},
\]
\[
y=\frac{c-a}{c-d}\cdot\frac{x-d}{x-a},
\]
\[
1-y=\frac{b-c}{c-b}\cdot\frac{c-a}{c-a},
\]
\[
1-k'y=\frac{b-c}{c-b}\cdot\frac{c-a}{c-a},
\]
and, thence,
\[
\sqrt{y\cdot 1-y\cdot 1-k'y}=\frac{(a-c)\sqrt{X}}{(c-a)\sqrt{(x-a)\cdot(x-a)}},
\]
where $\sqrt{\frac{(c-d)}{(c-d)}}$ is taken to be positive, and the signs of $\sqrt{X}$ is fixed as above. For $y$ between $0$ and $1$, $y=\frac{1}{x^2}$, $1-y$, $1-k'y$ will be positive, and $\sqrt{y}$ being between $1$ and $\frac{1}{x^2}$, $1-y$ will be negative, but $1-k'y$ will be positive, viz.\ the radical is then taken as a positive multiple of $i$.

We have moreover
\[
dy=\frac{c-a}{c-b}\cdot\frac{dx}{(x-a)^2},
\]
and therefore
\[
\frac{dy}{\sqrt{y\cdot 1-y\cdot 1-k'y}}=\sqrt{(c-d)\cdot(c-a)}\cdot\frac{dx}{\sqrt{X}},
\]
where $\sqrt{(c-d)\cdot(c-a)}$ is positive; or, say,
\[
\int\frac{dy}{\sqrt{y\cdot 1-y\cdot 1-k'y}}=\sqrt{(c-d)\cdot(c-a)}\cdot\int\frac{dx}{\sqrt{X}}.
\]

\paperref{\hfill $6-2$}

\newpage
% ============================================================
% PDF page 68 (printed 44)
% ============================================================
Hence, writing $y=\sin^2 u$, we have
\[
\Pi u=\sqrt{(c-d)\cdot(c-a)}\,\int_a^b\frac{dx}{\sqrt{X}},
\]
and it is to be further noticed that to
\[
x=a,\ b,\ c,\ d,
\]
correspond
\[
\Pi u=0,\ \tfrac{1}{2}K,\ \tfrac{1}{2}K+\tfrac{1}{2}K'.
\]
or we may say
\[
u=0,\ K,\ K+iK',\ 2K+iK'.
\]
Writing for shortness
\[
\frac{1}{\Pi}=\sqrt{(c-d)(c-a)}=u,
\]
we have
\[
\sin u=\int_a^b\frac{dx}{\sqrt{X}},
\]
and moreover
\[
\sin(K+iK')=\int_a^b\frac{dx}{\sqrt{X}},
\]
\[
n(K+iK')=\int_b^c\frac{dx}{\sqrt{X}},
\]
\[
n(2K+iK')=\int_c^d\frac{dx}{\sqrt{X}},
\]
or if for a moment we write
\[
\int_a^b\frac{dx}{\sqrt{X}}=A,\ \&c.
\]
then these equations are
\begin{align*}
nK&=A,\\
n(K+iK')&=A+B,\\
n(2K+iK')&=A+B+C.
\end{align*}
Hence $B+C=2A=B$, or $A=B-C$, or $B-C+D=0$, that is,
\[
\int_a^b\frac{dx}{\sqrt{X}}=\int_b^c\frac{dx}{\sqrt{X}},
\]
where observe as before that $x=a$ to $x=b$, or $x=c$ to $x=d$, $X$ is positive, and the radical $\sqrt{X}$ is taken to be positive.

We have also
\begin{align*}
nK&=A=\int_a^b\frac{dx}{\sqrt{X}},\\
niK'&=C-B=\int_a^b\frac{dx}{\sqrt{X}}.
\end{align*}

\newpage
% ============================================================
% PDF page 69 (printed 45)
% ============================================================
when, as before, from $b$ to $c$, $X$ is negative, and the sign of the radical is such that $\sqrt{X}$ is a positive multiple of $i$; the last formula may be more conveniently written
\[
n iK'=\int_b^c\frac{dx}{\sqrt{-X}},
\]
where, here $b$ to $c$, $-X$ is positive, and $\sqrt{(-X)}$ is to be taken as positive.

Collecting the results, we have
\[
\int_a^b\frac{dx}{\sqrt{X}}=\frac{nK}{\sqrt{(c-d)(c-a)}}, \ b\le a\le b,\ \dots b\le b\le b\le a,
\]
and also
\[
k=\frac{c-b}{a-b}\cdot\frac{d-a}{c-a},
\]
and thence conversely
\[
k'=\frac{a+d-b-c+(b-c-a+d)\,\sin^2 n}{a+b-c-d+(b-c+a-d)\,\sin^2 n},
\]
or, what is the same thing,
\begin{align*}
nu^2&=\frac{b-d\,\cdot\,x-a}{b-a\,\cdot\,x-d},\\
nu^2&=\frac{c-d\,\cdot\,x-a}{c-a\,\cdot\,x-d},\\
nu^2&=\frac{d-a\,\cdot\,x-b}{a-b\,\cdot\,x-d},\\
nu^2&=\frac{b-d\,\cdot\,x-a}{b-a\,\cdot\,x-d},
\end{align*}
where, in place of the elliptic functions $nu$ we substitute their $\vartheta$-values, it will be recollected that $\Omega$, the parameter of the $\vartheta$-functions, has the value
\[
\frac{K'}{K}\pi=\left(=\frac{\Omega'}{\Omega}\pi\right)=\pi\frac{\int_a^b\frac{dx}{\sqrt{-X}}}{\int_a^b\frac{dx}{\sqrt{X}}},
\]
and, so before,
\[
K=\int_a^b\frac{dx}{\sqrt{X}},
\]
Hence, finally, $a$, $b$, $K$, $K'$, $\Omega$ denoting given functions of $a$, $b$, $c$, $d$, as above
\[
\int_a^b\frac{dx}{\sqrt{X}}=nu,
\]
we have severally
\begin{align*}
\frac{b-d}{c-d}\cdot\frac{a-c}{c-a}&=\frac{1}{k^2}\,\frac{\vartheta_1^2 u}{\vartheta_2^2 u},&\frac{\vartheta_3^2 u}{\vartheta_2^2 u},&\\
\frac{d-a}{d-b}\cdot\frac{b-c}{a-c}&=\frac{k'^2}{k}\,\frac{\vartheta_2^2 u}{\vartheta_3^2 u},& \frac{\vartheta_1^2 u}{\vartheta_2^2 u}+\tfrac{1}{2}K',\\
\frac{c-a}{d-a}\cdot\frac{b-d}{b-c}&=k'^2\,\frac{\vartheta_4^2 u}{\vartheta_3^2 u},&\\
\end{align*}
which are the formulae in question.

\newpage
% ============================================================
% PDF page 70 (printed 46)
% ============================================================
The problem is to obtain these (and this in the more general case where $a$, $b$, $c$, $d$ have any given (real or imaginary) values directly from the assumed equation
\[
\int_a^b\frac{dx}{\sqrt{X}}=nu,
\]
and from the foregoing definition of the function $\vartheta u$.

It may be recalled that the function $\vartheta u$ is a doubly infinite product
\[
\vartheta u=\Pi\Pi\left[1+\frac{u}{m\pi+(n+\frac{1}{2})i\Pi}\right]
\]
as and a positive or negative integers from $-\infty$ to $+\infty$; I purposely omit all further explanations as to limits, or what is the same thing,
\[
\vartheta\frac{\pi}{2K}=\Pi\Pi\left[1+\frac{u}{2mK+(2n+1)iK'}\right]
\]
and consequently that, disregarding constant and exponential factors, the foregoing expressions of
\[
\frac{b-d}{c-d}\cdot\frac{a-c}{c-a},\quad \frac{d-a}{d-b}\cdot\frac{b-c}{a-c},\quad \frac{c-a}{d-a}\cdot\frac{b-d}{b-c},
\]
are the squares of the expressions
\[
\frac{X}{Y},\ \frac{Y}{Z},\ \frac{Z}{W},\quad \text{where } X,\ Y,\ Z,\ W \text{ are respectively of the form}
\]
\[
n\Pi\Pi\left[1+\frac{u}{(2n+1)iK'}\right],\quad \Pi\Pi\left[1+\frac{u}{(2m+1)iK'}\right],
\]
\[
\Pi\Pi\left[1+\frac{u}{2mK+iK'}\right],\quad \Pi\Pi\left[1+\frac{u}{(2m+1)K+iK'}\right],
\]
where $(m, n=2mK+2niK', \&c.$, the $smha$ over the $m$ or the $n$ denotes that the $2m$ or the $2n$ as the case may be) is to be changed into $2m+1$ or $2n+1$. But this is a transcription which has apparently no application to the $\vartheta$-functions of more than one variable.

\newpage
% ============================================================
% PAPER 717: ON THE TRIPLE THETA-FUNCTIONS
% PDF p71, printed 47
% ============================================================
\section*{717.}
\subsection*{ON THE TRIPLE THETA-FUNCTIONS.}

\begin{center}
\small [From the \textit{Messenger of Mathematics}, vol.~VII. (1878), pp.~48---50.]
\end{center}

\medskip

\textsc{As} a specimen of mathematical notation, viz.\ of the notation which appears to me the easiest to read, and also to print, I give the definition and demonstration of the fundamental properties of the triple theta-functions.

\textit{Definition.}
\[
\vartheta(U,\ V,\ W)=\Sigma\exp\Theta,
\]
where
\[
\Theta=(A,\ B,\ C,\ F,\ G,\ H)(l,\ m,\ n)^2+(U,\ V,\ W)(l,\ m,\ n).
\]
$l$, denoting the sum in regard to all positive and negative integer values from $-\infty$ to $+\infty$ (zero included) of $l$, $m$, $n$ respectively.

$\vartheta(U,\ V,\ W)$ is considered as a function of the arguments $(U,\ V,\ W)$, and it depends also on the parameters $(A,\ B,\ C,\ F,\ G,\ H)$.

\textit{First Property.} $\vartheta(U,\ V,\ W)=0$, for
\[
U=\tfrac{1}{2}\pi i+(A,\ H,\ G)(\alpha,\ \beta,\ \gamma),
\]
\[
V=\tfrac{1}{2}\pi i+(H,\ B,\ F)(\alpha,\ \beta,\ \gamma),
\]
\[
W=\tfrac{1}{2}\pi i+(G,\ F,\ C)(\alpha,\ \beta,\ \gamma),
\]
$\alpha$, $\beta$, $\gamma$, $\delta$, $\epsilon$, $\zeta$ being any positive or negative integer numbers, such that as $\alpha+\beta+\gamma=$ odd number.

\textit{Demonstration.} It is only necessary to show that to each term of $\vartheta$ there corresponds a second term, such that the inclusion of the two exponentials differ by an odd multiple of $\pi i$.

\newpage
% ============================================================
% PDF page 72 (printed 48)
% ============================================================
Taking $l$, $m$, $n$ the integers which belong to the one term, those belonging to the other term are
\[
-(l+a),\ -(m+\beta),\ -(n+\gamma).
\]
(where observe that one at least of the numbers $\alpha$, $\beta$, $\gamma$ being odd, this system of values is not in any case identical with $l$, $m$, $n$). The two exponents are
\[
\Theta_i=(A,\ B,\ C,\ F,\ G,\ H)(l,\ m,\ n)^2+(U,\ V,\ W)(l,\ m,\ n),
\]
and
\[
\Theta'_i=(A,\ B,\ C,\ F,\ G,\ H)(l+a,\ m+\beta,\ n+\gamma)^2-(U,\ V,\ W)(l+\alpha,\ m+\beta,\ n+\gamma),
\]
viz.\ the value of $\Theta'$ is
\[
=(A,\ B,\ C,\ F,\ G,\ H)(l,\ m,\ n)^2+2(A,\ H,\ G)(l\alpha,\ \beta,\ \gamma)+2(H,\ B,\ F)(m\alpha,\ \beta,\ \gamma)+2(G,\ F,\ C)(n\alpha,\ \beta,\ \gamma)
\]
\[
+(A,\ B,\ C,\ F,\ G,\ H)(a,\ \beta,\ \gamma)^2-(U,\ V,\ W)(l,\ m,\ n)-(U,\ V,\ W)(\alpha,\ \beta,\ \gamma),
\]
and we then have
\begin{align*}
\Theta'-\Theta&=-2(A,\ B,\ C,\ F,\ G,\ H)(l,\ m,\ n)^2+2(A,\ H,\ G)(\alpha,\ \beta,\ \gamma)l+2(H,\ B,\ F)(\alpha,\ \beta,\ \gamma)m\\
&\quad+2(G,\ F,\ C)(\alpha,\ \beta,\ \gamma)n+(A,\ B,\ C,\ F,\ G,\ H)(a,\ \beta,\ \gamma)^2\\
&\quad-2(U,\ V,\ W)(l,\ m,\ n)-(U,\ V,\ W)(\alpha,\ \beta,\ \gamma).
\end{align*}
which proves the theorem.

As to the notation, remark that, after $(A,\ B,\ C,\ F,\ G,\ H)$ has been written out in full, we may instead of $((A,\ B,\ C,\ F,\ G,\ H)(\alpha,\ \beta,\ \gamma)^2$ write
\[
(\alpha,\ \beta,\ \gamma)^2,
\]
and that we may use the like abbreviation
\[
(A,\ \dots)(\alpha,\ \beta,\ \gamma)\text{ for } (A,\ H,\ G)(\alpha,\ \beta,\ \gamma),\ \&c.,
\]
\[
(H,\ \dots)(\alpha,\ \beta,\ \gamma)\text{ for } (H,\ B,\ F)(\alpha,\ \beta,\ \gamma),\ \&c.,
\]
\[
(G,\ \dots)(\alpha,\ \beta,\ \gamma)\text{ for } (G,\ F,\ C)(\alpha,\ \beta,\ \gamma),\ \&c.
\]
These are not only abbreviations, but they make the formula actually clearer, as bringing them into a similar compact; and I accordingly use them in the demonstration which follows.

\newpage
% ============================================================
% PDF page 73 (printed 49)
% ============================================================
\textit{Second Property.} If $U$, $V$, $W$, denote
\begin{align*}
U'&=\pi i+(A,\ H,\ G)(\alpha,\ \beta,\ \gamma),\\
V'&=\pi i+(H,\ B,\ F)(\alpha,\ \beta,\ \gamma),\\
W'&=\pi i+(G,\ F,\ C)(\alpha,\ \beta,\ \gamma),
\end{align*}
respectively, where $\alpha$, $\beta$, $\gamma$, $\delta$, $\epsilon$, $\zeta$ are any positive or negative integers (zero values admissible), then
\[
\vartheta(U+u,\ V+v,\ W+w)=\exp\left[-2(\alpha U+\beta V+\gamma W)-(\alpha,\ \beta,\ \gamma)^2\right]\,\vartheta(U,\ V,\ W),
\]
or, say
\[
=\exp\left[-(A,\ \dots)(\alpha,\ \beta,\ \gamma)^2-2(U,\ V,\ W)(\alpha,\ \beta,\ \gamma)-(\alpha,\ \beta,\ \gamma)^2\right]\vartheta(U,\ V,\ W).
\]

\textit{Demonstration.} Writing $\vartheta(U,\ V,\ W)=\Sigma\exp\Theta$, then in the expression of $\Theta$, we may in place of $l$, $m$, $n$ write $-l-\alpha$, $-m-\beta$, $-n-\gamma$; we thus obtain
\[
\Theta_l=(A,\ \dots)(l,\ m,\ n)^2+2(A,\ \dots)(l,\ m,\ n)(a,\ \beta,\ \gamma)+(A,\ \dots)(\alpha,\ \beta,\ \gamma)^2
\]
\[
+(U,\ V,\ W)(l,\ m,\ n)+(U,\ V,\ W)(\alpha,\ \beta,\ \gamma),
\]
which is
\[
=(A,\ \dots)(l,\ m,\ n)^2
\]
\[
+2(\alpha,\ \beta,\ \gamma)(\alpha U+\beta V+\gamma W)-2(A,\ \dots)(\alpha,\ \beta,\ \gamma)+\Sigma'(A,\ \dots)(\alpha,\ \beta,\ \gamma)^2-\Sigma a^2(A,\ \dots)(\alpha,\ \beta,\ \gamma)
\]
\[
-2(\alpha,\ \beta,\ \gamma)^2+2\alpha,\ \beta,\ \gamma)+(A,\ \dots)(\alpha,\ \beta,\ \gamma)^2-2(A,\ \dots)(\alpha,\ \beta,\ \gamma)
\]
\[
+(U,\ V,\ W)(l,\ m,\ n)+(U,\ V,\ W)(\alpha,\ \beta,\ \gamma),
\]
which is
\[
=(A,\ \dots)(l,\ m,\ n)^2
\]
\[
+\Sigma'(A,\ \dots)(\alpha,\ \beta,\ \gamma)^2+\Sigma a^2(A,\ \dots)(\alpha,\ \beta,\ \gamma)+\Sigma\alpha\beta(A,\ \dots)(\alpha,\ \beta,\ \gamma)^2
\]
\[
+(U,\ V,\ W)(l,\ m,\ n)+\Sigma(A,\ \dots)(\alpha,\ \beta,\ \gamma)
\]
Hence, rejecting the last line, which (as on can readily verify) has the exponential multiplied as the case may be, $\vartheta(U,\ V,\ W)$ as $\vartheta(U,\ V,\ W)$ multiplied by the factor
\[
\exp\left[-(A,\ \dots)(\alpha,\ \beta,\ \gamma)^2-2(U,\ V,\ W)(\alpha,\ \beta,\ \gamma)-(\alpha,\ \beta,\ \gamma)^2\right],
\]
which is the theorem in question.

In many cases a formula, which belongs to an indefinite number of $k$ letters, is most easily intelligible when written out for three letters, but it is sometimes convenient to speak of the letters $l$, $m$, $\dots$ or even the letter $l$, and to write out the formula accordingly.

\paperref{C.~XI.\hfill 7}

\newpage
% ============================================================
% PDF page 74 (printed 50)
% ============================================================
In fact, writing herein $\alpha=\gamma=0$, that is, $a=\gamma=0$, the right-hand side becomes $a^2$ and the arts on the left-hand side are $a=B$, $A=B$, $a=A=I$, which suppose any four sum the sum of which is $=0$.

Writing in the last-mentioned equation $\alpha$, $\beta$, $\gamma$ for the $a$'s of $A$, $B$, $\gamma$, $\delta$ respectively, the equation becomes
\begin{align*}
P&=\beta P,\\
Q&=1-a^2+\beta^2+By^2,\\
R&=1-By^2-2\alpha+By^2,\\
S&=1-By^2.
\end{align*}
$L$, $M$, $N$ as the like functions $\alpha$, $\beta$, $\gamma$, $\delta$ respectively, $A$, $B$, $C$, $D$ being the same as before. And in fact
\begin{align*}
-L\,Q&=-y^2(\alpha^2+\alpha^2)+\beta^2+(\beta^2-\alpha^2)+B(y^2-x^2)+B(\beta-\alpha)+\dots,\\
&=\frac{1}{B}\left[\beta^2-(\beta-\alpha)^2+By^2-By^2\right]+B(y^2-x^2)+\dots,\\
&=\frac{1}{B}\left[\beta^2(x^2-y^2)+B(y^2-x^2)+B^2(y^2-x^2)\right],
\end{align*}
that is,
\[
-LP\,Q=-y^2(\beta^2-\alpha^2)+\beta^2-(\beta-\alpha)+\dots,
\]
\[
+(B+\alpha-\beta)\beta(y^2-x^2)+\beta-B^2(y^2-x^2)+\dots,
\]
or omitting the factor $\beta$, that is,
\[
-LP\,Q=-y^2(\beta^2-\alpha^2)+\beta^2-y^2(x^2-\alpha^2)+\dots+\beta(y^2-x^2)-By^2(x^2-\alpha^2),
\]
It is easy to verify the form of the orders 0, 1, 2, 3, and 4 in $x$, $y$, $z$, $w$ separately destroy each other, for instance, for the terms of the order 4, the value of $\Sigma$ is then
\[
=A^2+\Sigma\,y^2(\beta^2-\alpha^2)+y(x^2-\alpha^2)+B^2(x^2-\alpha^2)+(\Sigma\,\beta^2)(x^2-\alpha^2)+\dots,
\]
$\frac{1}{B}\left[\Sigma\,\beta\,a^2(x^2+y^2)+\Sigma\,a^2\,B^2(x^2-y^2)+B^2\,a^2(x^2-y^2)\right],$
It is easy to verify that the terms of the orders 0, 1, 2, 3, and 4 in $x$, $y$, $z$, $w$ separately destroy any sum sum of $x$ which is
\[
-LP^2+yQ^2+yP^2+xR^2+yS^2+(LP+yQ+xR+yS+yT)
\]
\[
+(D^2-1-LP^2)(yP^2+yP^2-yQ^2-yQ^2)+2(LP+yQ+xR+yS+yT)(LP+yQ+xR+yS+yT)=0,
\]
or, omitting the factor $\beta$, that is,
\[
-(LP+yQ^2+xR+yS+yT)+\Sigma(LP+yQ+xR+yS+yT)(LP+yQ+xR+yS+yT)+\dots=0,
\]
The theorem in the original form was obtained by me as follows: using the elliptic coordinates $p$, $q$, $r$, such that
\[
\frac{p^2}{a+p}+\frac{q^2}{b+p}+\frac{r^2}{c+p}=1,
\]

\newpage
% ============================================================
% PDF page 75 (printed 51)
% ============================================================
or, what is the same thing,
\[
-Bp\rho^2=a+p\,\cdot\,b+p\,\cdot\,c+p,
\]
\[
-y\eta y=b+q\,\cdot\,a+q\,\cdot\,c+q,
\]
\[
-z\zeta z=c+r\,\cdot\,a+r\,\cdot\,b+r,
\]
where $A$, $B$ denote $b-c$, $c-a$, $a-b$ respectively; then, treating $r$ as a constant, the coordinates $x$, $y$, $z$ will belong to a point on the ellipsoid
\[
\frac{x^2}{a+r}+\frac{y^2}{b+r}+\frac{z^2}{c+r}=1,
\]
and the differential equation of the right lines upon this surface is
\[
\frac{dp}{\sqrt{a+p\,\cdot\,b+p\,\cdot\,c+p}}+\frac{dq}{\sqrt{a+q\,\cdot\,b+q\,\cdot\,c+q}}=0.
\]
Take $a$, $b$, $c$, the coordinates of a point on the surface, and $p$, $q$, the corresponding values of $p$, $q$, so that
\[
-Bp_0\rho_0^2=a+p_0\cdot b+p_0\cdot c+p_0,
\]
\[
-Aq_0\eta_0=b+q_0\cdot a+q_0\cdot c+q_0,
\]
\[
-z\zeta_0=c+r\,\cdot\,a+r\,\cdot\,b+r,
\]
then the equation of the tangent plane at the point $(a, b, c, p_0, q_0)$ is
\[
\frac{x z_0}{a+r}+\frac{y y_0}{b+r}+\frac{z z_0}{c+r}=1,
\]
or, substituting for $x$, $y_0$, $z_0$ their values, we have
\[
=\beta p_0 \cdot ay+\cdot\,bxz+\,Az+\dots,
\]
and consequently the equation of the tangent plane through $(a, b, c, p_0, q_0)$ is the integral of the proposed differential equation.

Writing
\[
mz^2=A(a+p),\ mz^2=B(b+p),\ dx^2=C(c+p),
\]
the values of $A$, $B$, $C$, $S$ are determined; and, assuming for $q$, $r$, $p$, the like forms with the arguments $a$, $a_0$, $a_0$, the differential equation becomes $dz+dz$, having the
\paperref{\hfill $10-2$}

\newpage
% ============================================================
% PDF page 76 (printed 52)
% ============================================================
integral $u+x_0+v=u_0$; while the foregoing integral equation, on reducing the constant coefficients contained therein, takes the form
\[
-l^2 xx+mz_0\,mz+nz_0\,mz_0+nz_0\,nz+nz_0\,nz_0+\dots=0,
\]
\[
=
\frac{1}{2}lm\cdot mz_0\,mz_0,
\]
\[
=-l^2 nz \cdot mz_0\,mz_0,
\]
\[
=\frac{nz_0}{B^2-1-l^2}\cdot\frac{l^2}{l^2-l^2},
\]
via.\ this equation holds good if $a=u_0+v$, we have the theorem.

If, as above, $a=u+v+v$, $b=u-v+v$, $c=u+v-v$, where $r$ is regarded as constant, then the three products $\operatorname{snu\,nz\,nz}$, $\operatorname{cnu\,nz\,nz}$, $\operatorname{dnu\,nz\,nz}$, and equally so the three products $\operatorname{snz\,nz\,nz}$, $\operatorname{cnz\,nz\,nz}$, $\operatorname{dnz\,nz\,nz}$, are expressed by means of those products with the $u$, $v$, $w$ of the $u$, of the $v$, viz.\ for two surfaces it would be possible to obtain a real relation equation, which may be expressed in the squares of the $u$, $v$, $w$ of the $r$, of the $z$, all of whose internal coefficients are nominal, vis.\ the such would be obtained, in fact, a similar relation in the like form by giving to the parameter $r$ the advantage of being of a degree which is the half of the equation which involves only the squared functions.

Write $a$, $b$, $c$ or for $\operatorname{sn} u$, $\operatorname{cn} u$, $\operatorname{dn} u$, respectively; then, writing
\begin{align*}
A&=\sqrt{1-y^2\,\cdot\,1-l^2 y^2},& a&=\sqrt{1-l^2\,\cdot\,1-l^2 x^2},\\
B&=\sqrt{1-l^2\,\cdot\,1-l^2 y^2},& b&=\sqrt{1-x^2\,\cdot\,1-l^2 y^2},\\
C&=\sqrt{1-l^2\,\cdot\,1-l^2 y^2},& c&=\sqrt{1-l^2\,\cdot\,1-x^2-y^2},\\
D&=1-l^2 x^2y^2,& d&=1-l^2 z^2 x^2,
\end{align*}
we have
\[
\operatorname{sn}(z+v)=\frac{x A+y a}{D},
\]
that is,
\[
\frac{A+P}{D}=\frac{\sigma c}{\sigma a\,\cdot\,\sigma b},
\]
and consequently
\[
Dn=(A-x)(A+x),
\]
\[
PD=(P+x)(P-x);
\]
whence
\[
Dn=PD=(Aa-x)(Aa-x).
\]

\end{document}
